\section{How to read the formal evidence}
\label{sec:evidence-reading}
Paper X has a finite and formal core, a substantial arithmetic development, and a
terminal conditional assembly. These are not three names for one assertion.
The source register below gives the twenty numbered proof records in the main
manuscript. The checked endpoint concerns their recorded amended scope. Later
arithmetic prose can still require hypotheses which an application has not supplied.

A Lean declaration proves its type, including implicit parameters and typeclass
assumptions. A short proof can therefore depend on a long contract. Here the
terminal proof is only six lines: it constructs a conjunction and invokes two
previous theorems. The mathematical content is distributed across those theorems
and their dependencies. Reading only the final six lines would miss the source
family, all-depth naturality, inverse-limit closure, carrier presentations, and
constituent local comparisons.

\subsection{Three different records}
\begin{longtable}{@{}p{.20\linewidth}p{.75\linewidth}@{}}\toprule
Record & Interpretation\\\midrule
Block 27 and scalar/Mellin & Historical 17/20, with Claims 9.1, 9.2, and 11.1 still partial in those source revisions.\\
Conditional assembly & Recorded 20/20 under an explicitly amended contract. The target family and carrier summand are independent conclusions.\\
Later residual attempt & 19/20 and an unpromoted terminal claim under a stronger contract. Its unavailable pinned compiler is not a rejection of the earlier checked module.\\\bottomrule
\end{longtable}
The corresponding records are preserved independently. This edition does not
replace the later failed attempt by the earlier successful theorem, and it does
not demote the successful theorem merely because the later attempt was incomplete.

\subsection{What each kind of check can establish}
An exact Python calculation verifies the displayed finite instance or enumerated
range. A source hash verifies identity of bytes. A successful typeset build verifies
that the document can be generated. A saved compiler log reports a past
compilation. None of these is a fresh proof of every theorem in the collection.
A transitive axiom report also has a specific declaration set and environment.
Its scope is not enlarged by binding it into a longer monograph.

The present edition performs a new reader and release pass. It preserves the
formal sources and distinguishes inherited compiler evidence from checks actually
executed for this delivery. The dated final review and machine-readable verification
record list the execution results. They are the authoritative place for current
run counts, rather than numbers copied from an earlier checkpoint.

\clearpage
\section{The twenty source records}
\label{sec:register}
The numbered entries below are the main manuscript's twenty theorem, proposition,
and corollary records. Classical Inputs 3.1 and 4.1 are not counted as new proof
records. Definitions, examples, numerical certificates, and the conjectural
finite-seed-rank statement are also excluded from this count.

This is a reading concordance, not a claim that twenty lines of a table replace
twenty theorem-to-source audits. The earlier seventeen entries retain their
hypotheses. The final three arithmetic entries must be read through the amended
source contract explained in Sections~\ref{sec:horizontal}--\ref{sec:assembly}.
Links target named destinations in the companion monograph. Keep the three PDFs
in the same directory for relative file links to resolve.
\begin{longtable}{@{}p{.09\linewidth}p{.36\linewidth}p{.48\linewidth}@{}}\toprule
Record & Statement and source route & Hypothesis to retain\\\midrule\endfirsthead
\toprule Record & Statement and source route & Hypothesis to retain\\\midrule\endhead
\href{Paper_X_Traditional_Final.pdf#theorem.2.1}{2.1} & Residual partition. \par \small \code{Residual} & A supplied, single-valued reduction map.\\[5pt]
\href{Paper_X_Traditional_Final.pdf#theorem.2.2}{2.2} & No crossing under positive-depth perturbation. \par \small \code{Residual}, \code{SeedConsequences} & Root-preserving operations. The infinite statement assumes convergence.\\[5pt]
\href{Paper_X_Traditional_Final.pdf#theorem.2.3}{2.3} & Zeroth-order obstruction. \par \small \code{Residual.zeroth_order_obstruction} & Different roots cannot be reached by genuine positive-depth operations.\\[5pt]
\href{Paper_X_Traditional_Final.pdf#theorem.3.2}{3.2} & Formal inverse-limit reconstruction. \par \small \code{Reconstruction.represented_local_algebra_reconstruction} & Represented complete local problems with compatible finite maps.\\[5pt]
\href{Paper_X_Traditional_Final.pdf#theorem.3.3}{3.3} & Finite-branch compactness. \par \small \code{Compactness}, \code{RepresentationForests} & Finite nonempty levels and the stated total reduction maps.\\[5pt]
\href{Paper_X_Traditional_Final.pdf#theorem.3.4}{3.4} & Smooth formal branch count. \par \small \code{SmoothMaps}, \code{IntegralCharts} & An actual power-series chart and coefficient-preserving local maps.\\[5pt]
\href{Paper_X_Traditional_Final.pdf#theorem.4.2}{4.2} & Cohomology packet decomposition. \par \small \code{ComplexCohomology}, \code{PacketCohomology} & A finite complete orthogonal family of chain projectors.\\[5pt]
\href{Paper_X_Traditional_Final.pdf#theorem.5.3}{5.3} & Weighted smooth-chart count. \par \small \code{Counting}, \code{FormalCharts} & Independent coordinates with the stated valuation lower bounds.\\[5pt]
\href{Paper_X_Traditional_Final.pdf#theorem.6.2}{6.2} & Smooth seed zeta function. \par \small \code{Zeta}, \code{SmoothModels} & The smooth-chart count, as an infinite formal power series.\\[5pt]
\href{Paper_X_Traditional_Final.pdf#theorem.6.3}{6.3} & Weighted seed zeta function. \par \small \code{Zeta} & The stated ordered positive weights, with finite initial levels separated.\\[5pt]
\href{Paper_X_Traditional_Final.pdf#theorem.6.4}{6.4} & Singular power seed. \par \small \code{Singular}, \code{SingularZeta} & The scalar equation x to the e equals zero at finite precision.\\[5pt]
\href{Paper_X_Traditional_Final.pdf#theorem.6.6}{6.6} & Fractional singular dimension. \par \small \code{BranchingDimension}, \code{Singular} & A finite-count growth exponent, not the dimension of a limit space.\\[5pt]
\href{Paper_X_Traditional_Final.pdf#theorem.7.3}{7.3} & Directed minimum-seed theorem. \par \small \code{DirectedSeeds} & A finite supplied graph. Arithmetic admissibility of edges is separate.\\[5pt]
\href{Paper_X_Traditional_Final.pdf#theorem.8.2}{8.2} & Adelic finite-depth completion. \par \small \code{ClassCommonSpine}, \code{SystemCompletion} & Common data, finite-depth nonemptiness, and inverse-limit-closed conditions.\\[5pt]
\href{Paper_X_Traditional_Final.pdf#theorem.8.3}{8.3} & Residual companion obstruction. \par \small \code{SystemCompletion.system_empty_level_obstruction} & An empty required finite level rules out the specified realization.\\[5pt]
\href{Paper_X_Traditional_Final.pdf#theorem.9.1}{9.1} & Finite Frobenius separation. \par \small \code{Claim91ConstructedClosure} & Later constructed rational-base separator, with its literal compatible-family inputs.\\[5pt]
\href{Paper_X_Traditional_Final.pdf#theorem.9.2}{9.2} & Compatible-carrier summand completion. \par \small \code{StrictCompatibleCarrierCompletion} & Explicit copy presentations, natural multiplicities within capacity, and supplied local data.\\[5pt]
\href{Paper_X_Traditional_Final.pdf#theorem.10.1}{10.1} & Exact finite occupancy trace. \par \small \code{SpectralProjectors}, \code{JointOccupancy} & Commuting split semisimple operators. The positive-trace interpretation uses characteristic zero.\\[5pt]
\href{Paper_X_Traditional_Final.pdf#theorem.11.1}{11.1} & Residual-system completion criterion. \par \small \code{ConditionalOccupiedSystemCompletion} & The later amended two-conjunct contract, not universal residual existence.\\[5pt]
\href{Paper_X_Traditional_Final.pdf#theorem.11.2}{11.2} & Exact location of the remaining existence problem. \par \small \code{SpineExistenceCriterion}, \code{RequiredPacketOccupancy} & All required finite levels and all supplied packet occupancy premises.\\[5pt]
\bottomrule\end{longtable}
The monograph retains the general-base mathematical discussion. In contrast, the
printed constructed arithmetic separator and the terminal assembly use
\code{GQ}, the rational absolute Galois group. That literal base is part of the
formal statement. An arbitrary-number-field version requires its own attachment
and proof, not a change in a heading.

\clearpage
\section{Residual roots and reconstruction}
The fiber decomposition begins with a function, the reduction map. For an object
$x$ in a coefficient algebra $A$, its residual root is $\operatorname{red}_A(x)$.
Distinct fibers are disjoint because a function has one value at each input.
No topological or representation-theoretic theorem is needed for that partition.

The no-crossing theorem adds a hypothesis about the allowed operation. A
positive-depth correction is trivial after residual reduction. A finite sequence
of such corrections therefore stays in one fiber. An infinite sequence needs a
convergence assumption. The sequence $0,3,0,3,\ldots$ has one residue modulo three
but does not converge modulo nine. Residual invariance is weaker than being Cauchy.

\subsection{A point in an inverse limit is a coherent family}
For a represented complete deformation problem, compatible finite maps determine
one map into the inverse-limit coefficient ring. The argument does not assert
that every finite object can be extended. Nor does it produce a map when the
finite levels are empty. Reconstruction and existence remain distinct.

The finite-branch argument supplies existence when all levels are finite and
nonempty and the transition maps are the stated reductions. Surjectivity of every
transition map is unnecessary. Some nodes may die. The surviving branch is
obtained from nodes with descendants at arbitrarily large depths. Finiteness
prevents indefinitely moving the root to a new unrelated component.

Atlas V9 displays the infinite-level countermodel $X_m=\{n\ge m\}$ with inclusion
maps. Each level is nonempty, but a compatible family would have to be one integer
$n$ lying in every $X_m$. No such integer exists. Its infinite levels violate the
compactness hypotheses. This is a countermodel to dropping finiteness, not to the
finite-level theorem.

\subsection{The corrected digit formula}
For compatible canonical residues $0\le x_m<p^m$, the digit in position $m-1$ is
\[
 a_{m-1}=\left\lfloor\frac{x_m}{p^{m-1}}\right\rfloor\pmod p.
\]
Set $x_0=0$. Compatibility also makes
\[
 a_{m-1}=\frac{x_m-x_{m-1}}{p^{m-1}}\in\{0,\ldots,p-1\}.
\]
The earlier algorithm display omitted the floor. The executable implementation
already used integer division. For the base-five residues $4,19$, the new digit
is $3$. Dividing $19$ by $5$ in the rationals does not recover a digit. The new
regression checks compare the displayed formulas with the inherited executable
on exhaustive bounded towers, including noncanonical integer representatives
for the floor formula.

\clearpage
\section{Counting finite points without inventing limit points}
In a smooth chart $\mathcal O[[t_1,\ldots,t_d]]$, local maps into
$\mathcal O/\varpi^m$ send each parameter into the maximal ideal. There are
$q^{m-1}$ possibilities for each independent coordinate, giving
$q^{d(m-1)}$ maps. The word independent is justified by the power-series
presentation, not by the appearance of a grid in a figure.

Weighted coordinates with $t_i\in\varpi^{w_i}\mathcal O$ have
$\max(m-w_i,0)$ free residue digits. Thus
\[
 N_{\symbf{w}}(m)=q^{\sum_i\max(m-w_i,0)}.
\]
For weights $(1,2,4)$ over $\mathbb Z_2$, the counts begin
$1,2,8,32,256,2048$. The coordinate of weight four is not nontrivial at depth
four. It first has a free digit at depth five. The weight and first visible depth
are different quantities.

\subsection{The singular count is not a count of infinite branches}
For $x^e=0$ modulo $\varpi^m$, the condition is
$e\,v_\varpi(x)\ge m$. Its finite cardinality is
$q^{m-\lceil m/e\rceil}$. Splitting $m=ek+r$, $1\le r\le e$, gives
\[
 Z_e(T)=\frac{\sum_{r=1}^{e}q^{r-1}T^r}{1-q^{e-1}T^e}.
\]
This is an identity of formal power series. Its growth exponent $1-1/e$
measures the finite thickenings. It does not imply that the complete solution
space has a fractional geometric dimension. In the integral domain
$\mathcal O$, the only complete root of $x^e=0$ is zero. Equivalently, the
reduction at depth $m$ of a root at depth $em$ is zero. Every coherent infinite
branch is therefore the zero branch.

The smooth, weighted, and singular cases teach different tests. A complete chart
can have many genuine branches. A weighted chart delays their finite visibility.
A singular thickening can have many finite roots and only one complete branch.
The generating series must always be interpreted for the object it actually counts.

\clearpage
\section{Cohomology packets and directed transport}
A degreewise idempotent is not automatically a map of complexes. The required
identity is $d_rP_j^r=P_j^{r+1}d_r$. Together with orthogonality and a sum equal
to the identity, this makes the images into subcomplexes. The induced maps on
cycles and boundaries then descend to cohomology. Their inverse identities give
\[
 H^r(C^\bullet)\simeq\bigoplus_j H^r(P_jC^\bullet).
\]
The comparison with tangent or obstruction spaces of a particular deformation
problem is an additional identification. The packet theorem itself is a theorem
about a complex with supplied chain projectors.

Atlas V10 and V11 use a complex over $\mathbb F_5$ with dimensions $(2,4,2)$.
The valid packet family gives cohomology dimensions $(1,2,1)$. A modified
idempotent deletes the image of a retained source vector and fails the chain-map
identity. All its individual matrices remain idempotent. The missing condition
is compatibility with the differential, not idempotence.

\subsection{A transport graph is data before it is an optimization problem}
For a finite directed graph, contract strongly connected components. The resulting
graph is acyclic. Every generating seed set must meet each source component,
since no other component can reach into that source. Conversely, one vertex
from every source reaches all components. This proves the directed minimum-seed
formula.

The theorem optimizes a supplied graph. An arithmetic twist, induction, or
coefficient change is an edge only after its admissibility has been established.
Adding an unjustified arrow can lower the computed seed rank without proving
anything about the original arithmetic problem. Likewise, a finite graph theorem
does not settle the conjectural finite seed rank of an unrestricted arithmetic
transport category.

\subsection{Transport must respect every reduction square}
For roots of $x^2+1$ modulo powers of five, negation preserves the equation and
commutes with every reduction. Translation by one maps the root two to the root
three modulo five, but maps three to four, which is not a root. One successful
pair of endpoints does not define a natural operation on the whole tower.
The formal interface in the next section expresses exactly this distinction.

\clearpage
\section{Horizontal completion, hypothesis by hypothesis}
\label{sec:horizontal}
The module \code{HorizontalTransportCompletion} works with actual
change-of-basis classes. Its structure \code{AdmissibleClassTransport} has three
fields: a map at every precision, preservation of allowed classes, and
compatibility with reduction. Those fields are the data of a map of inverse towers.

A completed source representation $\rho$ first gives a coherent branch of its
finite classes. Applying the transport at each depth gives a coherent target
branch. The predecessor reconstruction theorem then supplies an actual target
representation. This is stronger than mapping one residual root, but it is still
conditional on the existence of the completed source representation.

\subsection{The local profile is not hidden in the word reconstruction}
\code{TransportedAlong} records, for every depth, equality between the reduced
target class and the transported source class. An arbitrary semantic property
of the complete representation does not follow from that equality alone.
The theorem therefore accepts \code{InverseLimitClosed} for the selected target
profile and finite-level conditions. Applying the closure theorem to the
reconstructed target establishes that profile.

For example, nontriviality is not a closed condition under the relevant limit.
The matrices
\[
 \rho_n(g)=\begin{pmatrix}1&3^n g\\0&1\end{pmatrix}
\]
are nontrivial representations of the additive integers, while they converge
pointwise to the trivial representation in the three-adic topology. A proof of
nontriviality for every approximation is not a proof of nontriviality for its limit.
This example does not decide the closedness of any particular arithmetic local
condition. That condition needs its own theorem.

\subsection{What changes when coefficient places are indexed}
The indexed theorem takes a source member at every $\lambda$. It first proves
an existence statement at each place, then uses choice to select a family of
target members. The source code explicitly verifies three outputs at every place:
finite-depth admissibility, the relation \code{TransportedAlong}, and membership
in the target profile.

This is an indexed family statement. Any interpretation as a compatible
arithmetic system depends on the shared data encoded by the conditions and
profiles. Merely placing unrelated local objects in one function indexed by
$\lambda$ is not a new compatibility theorem. The abstract horizontal module is
general in the source group, but the terminal assembly printed here specializes
that group to \code{GQ}.

\clearpage
\section{From arithmetic occurrence to a finite separator}
\label{sec:separator}
The constructed Claim 9.1 module has two declarations. The first provides its
\code{ExistentialArithmeticCoverage} proposition by calling the constructed
Chebotarev bridge. The second applies the arithmetic system-separation theorem
to that proposition. In this later source there is no free external coverage
parameter in the final separator signature. The earlier scalar/Mellin checkpoint
had retained that parameter. Both sources remain in the archive.

\subsection{What the separator accepts}
The family consists of finite-dimensional irreducible representations, compatible
character values in a common field, continuous characters at each realization,
and pairwise nonisomorphism in the stated memberwise sense. At each coefficient
place the bad-prime set is finite. The caller selects a distinguished place and
may forbid an additional finite set of rational primes.

The result selects distinct allowed rational primes. The common sample matrix
has nonzero determinant, and that remains true after each coefficient-field
embedding. Matrix entries identify with the representation characters wherever
the required good-prime condition holds. The terminal wrapper then converts this
matrix identity into the literal trace determinant needed by the carrier theorem.

Distinct primes alone would not suffice. Samples could repeat one character row
and leave the matrix singular. Nonisomorphism, arithmetic occurrence, and the
finite-dimensional separation argument do different work.

\subsection{The finite character-table model}
With columns trivial, sign, and standard, use
\[
 C=\begin{pmatrix}1&1&2\\1&-1&0\\1&1&-1\end{pmatrix},
 \qquad \det C=6,
 \qquad C\begin{pmatrix}2\\0\\1\end{pmatrix}
       =\begin{pmatrix}4\\2\\1\end{pmatrix}.
\]
The table explains multiplicity recovery after separating samples exist.
It does not produce rational primes with prescribed Frobenius behavior.
The arithmetic occurrence construction is represented by the preserved source
and its inherited verification record, not by this finite example.

\subsection{Exact recovery can still lose finite precision}
Over characteristic zero, $C$ is invertible. Modulo eight,
$\delta m=(4,4,0)$ is nonzero but $C\delta m=(8,0,8)$ vanishes. The denominator
valuation consumes one binary digit. Traces modulo sixteen distinguish this
particular ambiguity. The example separates invertibility over a field from
invertibility over a residue ring. Atlas V17 gives the complete small model.

\clearpage
\section{Carrier recovery is a construction inside supplied data}
\label{sec:carrier}
Let the supplied carrier be equivariantly presented as
$\bigoplus_i W_i^{\oplus c_i}$, and let the chosen multiplicities satisfy
$m_i\le c_i$. The copy inclusion and projection give maps $j$ and $r$ with
$r\circ j=\mathrm{id}$. Transporting these through the carrier presentation
produces an actual split summand. An arbitrary vector of traces does not provide
that presentation.

The initial subrepresentation is also supplied with an equivariant presentation
by the same selected copies. The nonsingular sample matrix proves that no other
natural multiplicity vector has those sample traces. Thus the theorem constructs
splittings and proves uniqueness within the declared constituent family. It does
not discover an unknown irreducible decomposition from raw observations.

\subsection{Three tests after solving a linear equation}
For capacities $(2,1,1)$, the vector $(2,0,1)$ is admissible. The trace vector
$(1,0,0)$ instead gives $(1/6,1/6,1/3)$, which is not integral. The vector
$(-1,1,1)$ is integral but negative. The vector $(3,0,1)$ exceeds capacity.
Each has a perfectly well-defined rational trace vector under the same matrix.
A nonsingular system alone does not turn its solution into a representation.

\subsection{What happens to local data}
The source constructs finite copies of supplied local packets. Characteristic
polynomials multiply with multiplicities. Relevant dimensions and traces add.
The supplied period comparisons similarly produce the selected finite-copy
comparison and weight-multiplicity formulas.

These identities do not attach Weil--Deligne or period data to an arbitrary
new representation. The data and comparisons are input. An extra semantic
condition is handled separately by \code{FiniteCopyClosed}. The terminal
\code{CarrierRecoveryAt} packages specific polynomial, local, and period
identities. It does not universally quantify over every possible local predicate.
The separate \code{selected_local_condition} theorem applies only after its
finite-copy-closure premise has been proved.

A fixed nontrivial determinant character illustrates the distinction. Two
one-dimensional constituents can each have determinant $\chi$, while their
direct sum has determinant $\chi^2$, not $\chi$. A determinant target must be
combined according to its actual product law. The word local does not make a
condition stable under arbitrary sums.

\clearpage
\section{Occupancy and the missing nonzero premise}
The identity $\operatorname{Tr}(P)=\operatorname{rank}(P)$ for a projector over a
characteristic-zero field is interpreted through the natural-number rank cast
into that field. Positivity means positivity of this integer, not an order on an
arbitrary coefficient field. In positive characteristic, a nonzero projector can
have field-valued trace zero. The rank remains the correct nonemptiness test.

For commuting split semisimple operators, products of their exact spectral
projectors detect joint eigenspaces. Pairwise occupancy need not imply joint
occupancy. Atlas V18 uses six binary labels for which each pair of the target
conditions $(0,1,1)$ is witnessed, while the complete triple is absent.
The product projector is zero although all pairwise tests pass.

\subsection{Why the terminal assembly permits an empty selection}
The multiplicities are natural numbers. The vector $m_i=0$ satisfies $m_i\le c_i$
for every carrier. Its selected copy space is zero, its splitting is valid, and
its polynomial products are empty products with value one. All of this is
consistent with the formal theorem.

To infer a nonzero summand, an application needs an index $i$ with $m_i>0$ and a
nonzero constituent $W_i$. Under the finite-dimensional field assumptions,
\[
 \dim\Bigl(\bigoplus_i W_i^{\oplus m_i}\Bigr)
       =\sum_i m_i\dim W_i>0
\]
then gives the desired conclusion. Dropping that premise would strengthen the
theorem to a false statement.

The indexed required-packet criterion is a different statement. It requires a
nonzero vector in every supplied required packet. A finite sample of packet
indices need not test them all. It also does not produce one common vector
belonging to several unrelated carriers. Quantifiers and carrier types must
remain visible in the application.

\clearpage
\section{The final conjunction and its exact logical strength}
\label{sec:assembly}
Write $A(\sigma)$ for the transported-family conclusion and $B(S)$ for carrier
recovery. The assembly has the form
\[
 (\exists\sigma\ A(\sigma))\ \wedge\ (\exists S\ B(S)).
\]
It does not assert an isomorphism between the two witnesses. In a finite logical
model, $A$ may hold only at zero and $B$ only at one. Both existential statements
are true, but no object satisfies both. This elementary countermodel explains
why conjoining proofs cannot create a missing comparison.

In the actual source, the separation is stronger than different names: the
transported family is represented by matrices over coefficient rings, whereas
carrier recovery is formulated through constituents and presentations over
coefficient fields. A comparison may need scalar extension, an equivariant
isomorphism, and compatibility with the local structures. There is no canonical
identification merely because both outputs are called a system or a packet.

\subsection{Read the six proof lines}
The terminal declaration first uses \code{constructor}. The left branch calls
\code{exists_transported_system_of_inverseLimitClosed}. The right branch calls
\code{carrier_recovery_from_arithmetic_separator}. The latter obtains a trace
separator and invokes the strict compatible-carrier theorem. The dependency
pattern is acyclic and the two branches share only the parameters they actually use.

The carrier statement is at an arbitrary distinguished coefficient place.
Repeating the theorem at all places requires the corresponding carrier and
initial-subrepresentation presentations there. Choice can select witnesses once
those pointwise existence statements have been proved. It cannot fill a place
at which their hypotheses are unavailable.

\subsection{A practical application contract}
An application should first identify its complete source members and the actual
finite-level transports. It must then prove the target profile's closure and
supply the constituent family, carrier presentations, and compatible sample data.
Only after applying the two construction branches should it attempt any desired
identification or nonzero conclusion. This order keeps construction, comparison,
and occupancy separate rather than using one as an unstated substitute for another.

\clearpage
\section{Source provenance and replay boundaries}
\label{sec:provenance}
The three direct dependency modules printed below are copied from the
\code{conditional_checked} archive tree. Their full leaf paths, nested-container
lineage, byte lengths, and SHA-256 values are recorded in
\code{SOURCE_EXCERPT_LINEAGE.json}. The terminal production module is also
retained as executable UTF-8 source. The printed listing is a view of that source,
not an alternate executable file.

\subsection{The compact preservation archive}
The object store deduplicates leaf payloads by SHA-256 and records each occurrence
in its original archive tree. Nested ZIP contents can be restored as directories.
This preserves their leaf bytes and lineage, not the original ZIP compression
streams or central-directory bytes. A container hash in the map identifies the
historical container. It is not a claim that a newly reconstructed ZIP has that hash.

From the release root, the following commands verify and restore the preserved
conditional tree. The output directory must be new or empty.
\begin{Verbatim}[fontsize=\small,breaklines=true]
python formal_preservation/restore.py --verify
python formal_preservation/restore.py --root conditional_checked --output restored_conditional
\end{Verbatim}
The original reproduction scripts and saved logs remain under the restored
lineage and readable records. They can contain original execution-workspace paths.
A fresh Lean replay requires the pinned toolchain, matching dependencies, and
corrected local paths. No such replay was performed for this reader edition.

\subsection{Interpreting the saved controls}
The historical verification reports two warnings-as-errors builds from separate
frozen-cache copies, matching corresponding proof objects, and a transitive axiom
audit of the exposed endpoint declarations. The allowed axioms are
\code{propext}, \code{Classical.choice}, and \code{Quot.sound}.
The false-proof and placeholder controls were rejected. Omitting closure or
separator premises was also rejected. The injected custom axiom compiled and was
then detected by the audit policy. It is inaccurate to say that Lean rejected it
as a malformed proof.

The pinned dependency cache and predecessor proof objects were trusted inputs to
that saved endpoint replay. Their existence in the archive does not mean this
publication rebuilt the full dependency graph. Source identity, saved compilation,
and a new cold build are separate claims.
